% ============================================================
%  Shared preamble for all builds (whole book and per-chapter).
%  Drivers: main.tex, main-ch01.tex, main-ch02.tex
% ============================================================

\usepackage[
  paperwidth=8.5in, paperheight=11in,
  twoside,
  inner=0.95in, outer=2.85in,
  top=1.0in, bottom=1.0in,
  marginparwidth=2.2in, marginparsep=0.2in,
  headsep=16pt
]{geometry}

\usepackage{annotated}
% Rudin's equation numbers are set with \tag, so several displays share an
% equation counter value and hyperref would mint the same destination name
% twice.  hypertexnames=false makes it use unique internal names instead.
\usepackage[colorlinks=true, allcolors=RuInk, hypertexnames=false]{hyperref}

\linespread{1.05}
\setlength{\parskip}{0pt}
\setlength{\parindent}{1.4em}

% Title page, parameterised by the chapter-coverage line and its subtitle.
\newcommand{\rutitlepage}[2]{%
\thispagestyle{empty}
\vspace*{1.5in}
\begin{center}
{\sffamily\Huge\bfseries\color{RuInk} The Annotated Rudin}\\[12pt]
{\sffamily\large\color{RuGrey} \emph{Principles of Mathematical Analysis}, Third Edition}\\[4pt]
{\sffamily\large\color{RuGrey} Walter Rudin}\\[36pt]
{\sffamily\normalsize\color{RuInk}\bfseries #1}\\[6pt]
{\sffamily\small\color{RuGrey} #2}
\end{center}
\vfill
\begin{center}
\begin{minipage}{0.8\textwidth}
\small\color{RuGrey}
Rudin's text appears here complete and unaltered --- his sentences, his
paragraph breaks, his numbering. Nothing has been paraphrased away. The
annotation is strictly additive: numbered markers in his prose key to notes in
the outer margin, and longer commentary sits in ruled boxes between his
paragraphs, where it can be skipped without breaking his argument.
\end{minipage}
\end{center}
\clearpage}
